\begin{trapbox}

	\begin{description}[style=nextline, leftmargin=2.8em, labelsep=0.5em, itemsep=0.35em]
		\item[\textbf{\textcolor{CTrap}{易错点一（漏除直线l2）}}]
			设直线系 $l_1+\lambda l_2=0$ 后，直接默认它能表示所有过交点直线，忽略了 $l_2$ 本身无法用有限 $\lambda$ 表示。
		\item[\textbf{\textcolor{CTrap}{正确理解（含l2检验）}}]
			使用单参数形式时，需额外判断 $l_2$ 是否满足题设条件，若满足应补上。
		\item[\textbf{\textcolor{CTrap}{错因分析（极限误解）}}]
			认为 $\lambda\to\infty$ 时得到 $l_2$，但 $\lambda$ 取有限实数永远不能恰好等于 $l_2$，故 $l_2$ 不在单参数系中。
	\end{description}

	\medskip
	\begin{description}[style=nextline, leftmargin=2.8em, labelsep=0.5em, itemsep=0.35em]
		\item[\textbf{\textcolor{CTrap}{易错点二（分母忽略）}}]
			在利用直线系求斜率或截距时，不讨论分母为零的情况，导致失去有效解。例如由 $k=-\frac{A_1+\lambda A_2}{B_1+\lambda B_2}$ 直接列方程，未考虑 $B_1+\lambda B_2=0$ 的情形。
		\item[\textbf{\textcolor{CTrap}{正确理解（分类讨论）}}]
			遇到含分式表达的参数关系时，应分类讨论分母为零和不为零两种情况，分别求解。
		\item[\textbf{\textcolor{CTrap}{错因分析（惯性约分）}}]
			学生习惯直接十字相乘或约去公因式，忽略了分母为零是方程成立的一种可能，造成失根。
	\end{description}

	\medskip
	\begin{description}[style=nextline, leftmargin=2.8em, labelsep=0.5em, itemsep=0.35em]
		\item[\textbf{\textcolor{CTrap}{易错点三（参数范围不清）}}]
			使用双参数形式 $m(A_1x+B_1y+C_1)+n(A_2x+B_2y+C_2)=0$ 时，忘记规定 $m,n$ 不同时为零，导致方程退化为恒等式或矛盾式。
		\item[\textbf{\textcolor{CTrap}{正确理解（非零约束）}}]
			双参数中 $m,n$ 不能同时为零；通常可设其中一个为 $1$（但不能覆盖形如 $m=0,n\neq0$ 的情形），最好明确 $m^2+n^2\neq0$。
		\item[\textbf{\textcolor{CTrap}{错因分析（参数缺失）}}]
			认为与单参数类似，忽略双参数系中 $m=0,n=1$ 代表 $l_2$，若设 $m=1$ 则漏掉 $l_2$，与易错点一相同。
	\end{description}

\end{trapbox}
